\documentclass[11pt,a4paper]{article}

% =========================================================================
% PAKET-PAKET UTAMA
% =========================================================================
\usepackage[utf8]{inputenc}
\usepackage[T1]{fontenc}
\usepackage{lmodern}
\usepackage[indonesian]{babel}
\usepackage{amsmath,amssymb,amsfonts,amsthm}
\usepackage{bm}              % Notasi tebal untuk vektor & matriks
\usepackage{geometry}
\geometry{top=2.5cm, bottom=2.5cm, left=2.5cm, right=2.5cm}
\usepackage{graphicx}
\usepackage{booktabs}        % Tabel profesional
\usepackage{tabularx}
\usepackage{microtype}
\usepackage{xcolor}
\usepackage{mdframed}        % Kotak highlight untuk definisi
\usepackage{caption}
\usepackage{subcaption}
\usepackage{float}
\usepackage{hyperref}
\hypersetup{
    colorlinks=true,
    linkcolor=blue!70!black,
    citecolor=blue!70!black,
    urlcolor=blue!70!black,
    pdfencoding=auto,
    psdextra
}
\pdfstringdefDisableCommands{%
  \def\boldsymbol#1{#1}%
  \def\bm#1{#1}%
  \def\mathbf#1{#1}%
}

% Makro sitasi fleksibel
\newcommand{\citep}[2][]{(Rasmussen \& Williams, 2006\if\relax\detokenize{#1}\relax\else, #1\fi)}
\newcommand{\citet}[2][]{Rasmussen \& Williams (2006\if\relax\detokenize{#1}\relax\else, #1\fi)}

% Path direktori gambar
\graphicspath{{figures/}}

% =========================================================================
% INFORMASI DOKUMEN
% =========================================================================
\title{
    \vspace{-1.2cm}
    \textbf{\Large BIMBINGAN TUGAS AKHIR \#4} \\[0.2em]
}
\author{}
\date{}

\begin{document}

\maketitle
\vspace{-1.2cm}

% =========================================================================
% 1. MARGINAL LIKELIHOOD (TYPE II MAXIMUM LIKELIHOOD / ML-II)
% =========================================================================
\section{Marginal Likelihood (Type II Maximum Likelihood / ML-II)}
\label{sec:marginal_likelihood}
\noindent \textit{\citep[Bab~5.4.1, hlm.~112--113, Persamaan~(5.7)--(5.8)]{rasmussen2006gaussian}}

\begin{equation}
p(\mathbf{y} \mid X, \boldsymbol{\theta}) = \int p(\mathbf{y} \mid \mathbf{f}) \, p(\mathbf{f} \mid X, \boldsymbol{\theta}) \, d\mathbf{f}
\end{equation}

\begin{equation}
\mathbf{y} \mid X, \boldsymbol{\theta} \sim \mathcal{N}\big(\mathbf{0}, K_y\big), \quad K_y = K + \sigma_n^2 I_N
\end{equation}

\begin{equation}
\log p(\mathbf{y} \mid X, \boldsymbol{\theta}) = -\frac{1}{2} \mathbf{y}^T K_y^{-1} \mathbf{y} - \frac{1}{2} \log |K_y| - \frac{N}{2} \log(2\pi)
\end{equation}

\begin{equation}
\operatorname{NLML}(\boldsymbol{\theta}) = \frac{1}{2} \mathbf{y}^T K_y^{-1} \mathbf{y} + \frac{1}{2} \log |K_y| + \frac{N}{2} \log(2\pi)
\end{equation}

% =========================================================================
% 2. MAXIMUM A POSTERIORI (MAP)
% =========================================================================
\section{Maximum A Posteriori (MAP)}
\label{sec:map}
\noindent \textit{\citep[Bab~5.2, hlm.~108--109, Persamaan~(5.2)]{rasmussen2006gaussian}}

\begin{equation}
p(\boldsymbol{\theta} \mid \mathbf{y}, X) = \frac{p(\mathbf{y} \mid X, \boldsymbol{\theta}) \, p(\boldsymbol{\theta})}{p(\mathbf{y} \mid X)}
\end{equation}

\begin{equation}
\mathcal{L}_{\text{MAP}}(\boldsymbol{\theta}) = \log p(\mathbf{y} \mid X, \boldsymbol{\theta}) + \log p(\boldsymbol{\theta})
\end{equation}

\begin{equation}
\boldsymbol{\theta}_{\text{MAP}} = \arg\max_{\boldsymbol{\theta}} \left[ -\frac{1}{2} \mathbf{y}^T K_y^{-1} \mathbf{y} - \frac{1}{2} \log |K_y| + \log p(\boldsymbol{\theta}) \right]
\end{equation}

% =========================================================================
% 3. CROSS-VALIDATION: LEAVE-ONE-OUT SQUARED ERROR (LOO-SE)
% =========================================================================
\section{Leave-One-Out Cross-Validation: Squared Error}
\label{sec:loo_se}
\noindent \textit{\citep[Bab~5.3 \& Bab~5.4.2, hlm.~111 \& hlm.~116--117, Persamaan~(5.5) \& (5.12)]{rasmussen2006gaussian}}

\begin{equation}
\mu_{-i} = y_i - \frac{[K_y^{-1}\mathbf{y}]_i}{[K_y^{-1}]_{ii}}
\end{equation}

\begin{equation}
L_{\text{LOO-SE}}(\boldsymbol{\theta}) = \sum_{i=1}^N (y_i - \mu_{-i})^2 = \sum_{i=1}^N \left( \frac{[K_y^{-1}\mathbf{y}]_i}{[K_y^{-1}]_{ii}} \right)^2
\end{equation}

% =========================================================================
% 4. CROSS-VALIDATION: LEAVE-ONE-OUT NEGATIVE LOG PSEUDO-LIKELIHOOD (LOO-NLPL)
% =========================================================================
\section{LOO Cross-Validation: Negative Log Pseudo-Likelihood}
\label{sec:loo_nlpl}
\noindent \textit{\citep[Bab~5.3 \& Bab~5.4.2, hlm.~111 \& hlm.~116--118, Persamaan~(5.6), (5.11) \& (5.13)]{rasmussen2006gaussian}}

\begin{equation}
\sigma_{-i}^2 = \frac{1}{[K_y^{-1}]_{ii}}
\end{equation}

\begin{align}
L_{\text{LOO-NLPL}}(\boldsymbol{\theta}) &= -\sum_{i=1}^N \log p(y_i \mid X, \mathbf{y}_{-i}, \boldsymbol{\theta}) \nonumber \\
&= \frac{1}{2} \sum_{i=1}^N \left( \log(2\pi \sigma_{-i}^2) + \frac{(y_i - \mu_{-i})^2}{\sigma_{-i}^2} \right) \nonumber \\
&= \frac{1}{2} \sum_{i=1}^N \left( -\log [K_y^{-1}]_{ii} + \frac{([K_y^{-1}\mathbf{y}]_i)^2}{[K_y^{-1}]_{ii}} + \log(2\pi) \right)
\end{align}

% =========================================================================
% DAFTAR PUSTAKA
% =========================================================================
\begin{thebibliography}{99}

\bibitem{rasmussen2006gaussian}
Rasmussen, C.~E., \& Williams, C.~K. (2006). \textit{Gaussian Processes for Machine Learning}. MIT Press, Cambridge, MA.

\end{thebibliography}

\end{document}
